\section{Transverse rigidity of bounded eigenfunctions}\label{sec:rigidity}

The coefficient estimates of Section~\ref{sec:energy} do not yet give
square summability of the horizontal derivatives.  We obtain that summability
from an identity involving the two modes at the lower end of an angular
Fourier tail.
The long averages used below adapt the construction for invariant functions
in \cite[Sections~4--5]{Companion}; the phase in the average and the endpoint
terms retain the spectral parameter throughout the argument.

\begin{proof}[Proof of Proposition~\ref{prop:rigidity}]
Write $H=\norm{f}_\infty$ and $A=X-i\lambda$.  The eigenfunction equation
implies $Af=0$ locally in distributions.  Lemma~\ref{lem:bounded-modes}
therefore gives
\begin{equation}\label{eq:rigidity-mode-bound}
 \norm{\nabla f_j}_2\le C_\lambda H
 \qquad(j\in\Z).
\end{equation}
Here and below constants may depend on the fixed triangle, but not on an
angular index or on the approximation parameters.

\paragraph{The two boundary modes of a Fourier tail.}
Fix $m\in\Z$.  Define the $L^2$ field and the two coefficient derivatives
\[
 U=\sum_{j\ge m}P_jf,\qquad
 b=\pb f_m,\qquad a=\pa f_{m-1},
\]
and put
\[
 B=b(x)e^{i(m-1)\theta}-a(x)e^{im\theta}.
\]
The factors in this expression respect the coefficient transition rules,
so $B$ is a global $L^2(SM)$ field.  The mode equations give
\begin{equation}\label{eq:tail-equation}
 AU=B\qquad\text{locally in distributions on }SM.
\end{equation}
For completeness, for $N\ge m+2$ the finite tail satisfies
\[
 A\sum_{j=m}^N P_jf
 =B-(\pb f_{N+1})e^{iN\theta}
      +(\pa f_N)e^{i(N+1)\theta}.
\]
The last two terms are bounded in $L^2$ by
\eqref{eq:rigidity-mode-bound} and converge weakly to zero: their angular
frequencies tend to infinity.  Since the finite tails converge strongly
to $U$ in $L^2$, this proves \eqref{eq:tail-equation}.

We will prove the endpoint identity
\begin{equation}\label{eq:endpoint-identity}
 \ip{B}{(P_{m-1}+P_m)Yf}=0.
\end{equation}
At this stage $Yf$ denotes a distribution; its two displayed modes are
in $L^2$ by \eqref{eq:modes} and \eqref{eq:rigidity-mode-bound}.
The difficulty is that $Yf$ itself is not yet known to be an admissible
test in \eqref{eq:tail-equation}.  We construct compactly supported
approximations whose error under $A$ becomes negligible in that pairing.

\paragraph{Long averages away from the vertices.}
Let $K$ be the constant in Lemma~\ref{lem:smoothing} and fix $L>K+4$.
For $0<s<1$ and sufficiently small $\eps>0$, set $T=s/\eps$.
Let $G_\eps$ consist of the regular states $z$ whose entire trajectory
segment $\{\Psi_tz:0\le t\le T\}$ lies over $\{r>L\eps\}$.
Then
\begin{equation}\label{eq:bad-segments}
 \nu(SM\setminus G_\eps)\le C(s+\eps^2).
\end{equation}
Indeed, $r$ is $1$-Lipschitz along a unit-speed geodesic.  If a segment
meets $\{r\le L\eps\}$, a time mesh of spacing at most $\eps$ detects
a visit to $\{r\le(L+1)\eps\}$.  There are at most $T/\eps+2$ mesh
points.  Measure preservation and the cone-area bound
$\nu\{r\le(L+1)\eps\}\le C\eps^2$ prove
\eqref{eq:bad-segments}.  The discarded trajectories that hit a vertex
form a null set.

Define
\begin{equation}\label{eq:segment-average}
 q_\eps=f\ind_{G_\eps},\qquad
 g_\eps=\frac1T\int_0^T e^{i\lambda t}
                      q_\eps\circ\Psi_{-t}\,dt.
\end{equation}
The integral may be taken in $L^2$; bounded measurable representatives
and Fubini give the same field.  The phase is chosen so that replacing
$q_\eps$ by $f$ in the integrand gives $f$.  Consequently,
\begin{equation}\label{eq:average-bounds}
 \norm{g_\eps}_\infty\le H,\qquad
 \norm{g_\eps-f}_2\le CH\sqrt{s+\eps^2},\qquad
 \norm{Ag_\eps}_2\le\frac{2H}{T}.
\end{equation}
The first bound follows directly from \eqref{eq:segment-average}; the
second follows by integrating the $L^2$ distance
$\norm{q_\eps-f}_2\le H\nu(G_\eps^c)^{1/2}$.
For the third we use the distributional identity
\begin{equation}\label{eq:average-endpoints}
 Ag_\eps=\frac1T
       \bigl(q_\eps-e^{i\lambda T}q_\eps\circ\Psi_{-T}\bigr).
\end{equation}
Here no derivative of $\ind_{G_\eps}$ is taken.  To verify the identity,
test against $\varphi\in C_c^\infty(SM)$ and change variables by the
flow.  Along every regular orbit,
\[
 \frac{d}{dt}\left(e^{i\lambda t}
                  \overline{\varphi(\Psi_tz)}\right)
 =e^{i\lambda t}\overline{(A\varphi)(\Psi_tz)}.
\]
Integrating this formula in $t$, and using the formal skew adjointness
of $A$, proves \eqref{eq:average-endpoints}.  The test and its horizontal
derivative are bounded, which also justifies Fubini in this calculation.

Both $g_\eps$ and $Ag_\eps$ vanish over $\{r<L\eps\}$.  In fact, for
each $0\le t\le T$ the support of
$q_\eps\circ\Psi_{-t}$ is contained, up to null sets, in
$\Psi_tG_\eps\subset\{r>L\eps\}$; this also applies to both endpoint
terms in \eqref{eq:average-endpoints}.
Apply the spatial smoothing of Lemma~\ref{lem:smoothing} and set
\[
 w_\eps=S_\eps g_\eps.
\]
Before its prescribed zero extension this field vanishes on
$\{K\eps<r<(L-1)\eps\}$.  Thus the extension is smooth in position,
has compact spatial support in $M$, and creates no boundary derivative.
Local commutation with $A$ gives, globally as an $L^2$ identity,
$Aw_\eps=S_\eps Ag_\eps$.  The smoothing estimates imply
\begin{equation}\label{eq:regularized-bounds}
 \begin{aligned}
 \norm{w_\eps}_\infty&\le H,
 &\norm{Aw_\eps}_2&\le\frac{CH}{T},\\
 \norm{\nabla w_\eps}_2&\le\frac{CH}{\eps},
 &\norm{YAw_\eps}_2&\le\frac{CH}{T\eps}=\frac{CH}{s}.
 \end{aligned}
\end{equation}
All higher spatial derivatives also belong to $L^2$ for fixed $\eps$.
The relation $T\eps=s$ bounds the discarded measure by $C(s+\eps^2)$
while keeping $YAw_\eps$ bounded in $L^2$ at fixed $s$. Its norm need
not tend to zero: weak convergence suffices for pairing against the fixed
tail $U$. We first let $\eps$ tend to zero at fixed $s$ to obtain exact
bounded local solutions. Only then can Lemma~\ref{lem:bounded-modes}
restore coefficient estimates independent of $s$ and allow the limit
$s\to0$.

\paragraph{The first limit: $\eps\to0$ at fixed $s$.}
The last bound in \eqref{eq:regularized-bounds} makes $YAw_\eps$ bounded
in $L^2$ at fixed $s$.  Since $Aw_\eps\to0$ in $L^2$, for each
$\varphi\in C_c^\infty(SM)$ we have
\[
 \ip{YAw_\eps}{\varphi}
 =-\ip{Aw_\eps}{Y\varphi}\longrightarrow0.
\]
Density of these tests in $L^2(SM)$ therefore yields
\begin{equation}\label{eq:weak-error}
 YAw_\eps\rightharpoonup0\quad\text{in }L^2(SM).
\end{equation}
We may now test \eqref{eq:tail-equation} against $Yw_\eps$ and commute
the horizontal fields:
\begin{equation}\label{eq:tail-pairing}
 \ip{B}{Yw_\eps}
 =-\ip{U}{AYw_\eps}
 =-\ip{U}{YAw_\eps}\longrightarrow0.
\end{equation}
This use of a test with only $L^2$ angular regularity is legitimate.
Indeed, the finite angular Fourier sums of $Yw_\eps$ are smooth compactly
supported tests.  They converge to $Yw_\eps$, together with their first
horizontal derivatives, in $L^2$, because $w_\eps$ has square-integrable
spatial derivatives of every order.  Apply
\eqref{eq:tail-equation} first to these finite sums and then pass to the
limit.  The equality $AYw_\eps=YAw_\eps$ follows from $[X,Y]=0$.

To identify the left side of \eqref{eq:tail-pairing} after taking a
limit, we also need coefficient-gradient bounds.  Lemma~\ref{lem:energy},
with $E=M$, and \eqref{eq:regularized-bounds} give
\begin{equation}\label{eq:fixed-s-mode-bound}
 \norm{\nabla(w_\eps)_j}_2^2
 \le CH^2(s^{-1}+\lambda^2)
 \qquad(j\in\Z).
\end{equation}
Choose a sequence $\eps\downarrow0$ along which $w_\eps$ converges weakly
in $L^2$ to a field $h_s$.  Then
\begin{equation}\label{eq:first-limit}
 \norm{h_s}_\infty\le H,
 \qquad Ah_s=0\text{ locally in distributions},
 \qquad \norm{h_s-f}_2\le CH\sqrt{s}.
\end{equation}
For the first assertion use weak closedness in $L^2$ of the convex set
$\{u:|u|\le H\ \text{a.e.}\}$; the second follows from
$Aw_\eps\to0$ in $L^2$.  For the third, write
\[
 w_\eps-f=S_\eps(g_\eps-f)+(S_\eps f-f)
\]
and use \eqref{eq:average-bounds}, uniform $L^2$ boundedness of
$S_\eps$, and $S_\eps f\to f$ in $L^2$.

Only the coefficients with indices $m-2,m-1,m,m+1$ contribute to the
two modes $(P_{m-1}+P_m)Yw_\eps$.  By
\eqref{eq:fixed-s-mode-bound}, a further subsequence makes their
gradients converge weakly in the corresponding global $L^2$ coefficient
spaces.  Testing in regular coordinate charts identifies these limits
as the gradients of the respective coefficients of $h_s$.
Since $B$ has only modes $m-1$ and $m$, \eqref{eq:tail-pairing} gives
\begin{equation}\label{eq:first-limit-pairing}
 \ip{B}{(P_{m-1}+P_m)Yh_s}=0.
\end{equation}

\paragraph{The second limit: $s\to0$.}
The bound \eqref{eq:fixed-s-mode-bound} deteriorates as $s\to0$.
The first limit has replaced the approximate equation for $w_\eps$
by the exact local equation $Ah_s=0$, while preserving boundedness.
Lemma~\ref{lem:bounded-modes} now gives the stronger estimate
\[
 \norm{\nabla(h_s)_j}_2\le C_\lambda H
 \qquad(j\in\Z),
\]
independently of $s$.  No flow-invariance assertion for $h_s$ is needed
here: the lemma requires only its boundedness and its local
distributional equation.

Choose any sequence $s\downarrow0$.  By \eqref{eq:first-limit},
$h_s\to f$ strongly in $L^2$.  After passing to a subsequence, the
gradients of the same four coefficients converge weakly in $L^2$;
local distributional testing identifies their limits as those of $f$.
Passing to the limit in \eqref{eq:first-limit-pairing} proves
\eqref{eq:endpoint-identity}.  The successive limits have thus removed
the approximation error without losing the coefficient derivatives
needed by the endpoint pairing.

\paragraph{The sharp horizontal-energy bound.}
By $Af=0$ and the mode formulas \eqref{eq:modes},
\[
 (Yf)_{m-1}=-\lambda f_{m-1}-2ib,\qquad
 (Yf)_m=\lambda f_m+2ia.
\]
Substitute these expressions into \eqref{eq:endpoint-identity}.
With inner products linear in the first argument, the result is
\[
 2i\bigl(\norm{b}_2^2+\norm{a}_2^2\bigr)
 =\lambda\bigl(\ip{b}{f_{m-1}}+\ip{a}{f_m}\bigr).
\]
Taking absolute values and applying Cauchy--Schwarz gives
\begin{equation}\label{eq:endpoint-energy}
 \norm{\pb f_m}_2^2+\norm{\pa f_{m-1}}_2^2
 \le\frac{\lambda^2}{4}
       \bigl(\norm{f_{m-1}}_2^2+\norm{f_m}_2^2\bigr).
\end{equation}
This remains valid when the left side is zero; otherwise divide the
Cauchy--Schwarz inequality by its square root.
Summing \eqref{eq:endpoint-energy} over all $m\in\Z$ yields
\begin{equation}\label{eq:sharp-energy}
 \sum_{j\in\Z}\norm{\nabla f_j}_2^2
 =2\sum_{j\in\Z}
     \bigl(\norm{\pa f_j}_2^2+\norm{\pb f_j}_2^2\bigr)
 \le\lambda^2\norm{f}_2^2.
\end{equation}
In particular the horizontal derivatives of the angular partial sums
converge in $L^2$; their limits are the weak derivatives of $f$, by
local testing.  Parseval and \eqref{eq:gradient} therefore give
\[
 \norm{Xf}_2^2+\norm{Yf}_2^2
 =\norm{\nabla f}_2^2
 \le\lambda^2\norm{f}_2^2.
\]
But $Xf=i\lambda f$, so the first term already equals the right side.
It follows that $Yf=0$ and equality holds throughout.
\end{proof}
